\documentclass[paper=a6,landscape,ngerman,twoside=true,fontsize=12pt]{scrartcl}
% ---------------------------------------------------------------------
% praeambel_app_standalone.tex
% ---------------------------------------------------------------------

% ---- Page shape ----
% ---- Display-Format für Smartphones (z. B. 90mm x 160mm) ----
\usepackage[
    paperwidth=105mm,
    paperheight=148mm,
    margin=5mm
]{geometry}

% font
\usepackage[sfdefault]{roboto}
\usepackage[T1]{fontenc}
\usepackage{eulervm}


% ---- Language / Encoding ----
\usepackage[utf8]{inputenc}
\usepackage[T1]{fontenc}
\usepackage[ngerman]{babel}   

% ---- Math ----
\usepackage{amsmath}          
\usepackage{amssymb}          
\usepackage{extarrows}        

% ---- Tables ----
\usepackage{array}            
\usepackage{tabularx}
\usepackage{diagbox}          

% Defined custom column type 'C' for fixed-width centered columns (e.g. NEWNEWNEW rule)
\newcolumntype{C}[1]{>{\centering\arraybackslash}p{#1}}

% ---- Graphics / Diagrams ----
\usepackage{graphicx}         
\usepackage{tikz}             
\usetikzlibrary{patterns}

% ---- Custom Math Commands ----
\newcommand{\R}{\mathbb{R}}                     
\newcommand{\D}{\mathbb{D}}                     
\newcommand{\abs}[1]{\left|#1\right|}           
\newcommand{\mal}{\cdot}                        
\newcommand{\std}{\sigma}                       

% ---- Card Environments ----
% ---- Card Environments (für Mobilgeräte optimiert) ----
\newenvironment{frage}{%
	\clearpage
	\vspace*{\fill}
	\begin{center}\huge
}{%
	\end{center}
	\vspace*{\fill}
}

\newenvironment{antwort}{%
	\clearpage
	\vspace*{\fill}
	\begin{center}\large
}{%
	\end{center}
	\vspace*{\fill}
}

\begin{document}

% category: Analysis
\begin{frage}
	Nullstellen einer Funktion?
	
\end{frage}
\begin{antwort}
	Ansatz: Funktion mit 0 gleichsetzen
	\[
	f(x)=0
	\]
\end{antwort}
\begin{frage}
	Nullstellen: Lineare Funktion
		
	z.B.: $f(x)=4x-12$
		
\end{frage}
\begin{antwort}
	Einfach umstellen:
	\begin{align*}
	4x-12&=0 \\
	4x&=12 \\
	x&=3
	\end{align*}
\end{antwort}
\begin{frage}
	Nullstellen: Quadratische Funktion
	\[
		f(x)=2x^2+6x-8
	\]
\end{frage}
\begin{antwort}
		\[
		2x^2+6x-8=0
		\]
	"`Mitternachtsformel"'
	\[
		x_{1,2}=\frac{-b\pm\sqrt{b^2-4ac}}{2a}
	\]
	\[
	x_{1,2}=\frac{-6\pm\sqrt{6^2-4\cdot2\cdot(-8)}}{2\cdot2 }
	\]
	\[
	x_{1}=1,\;x_{2}=-4
	\]

\end{antwort}
\begin{frage}
	Nullstellen: Funktion 3. Grades
\end{frage}
\begin{antwort}
	\normalsize\begin{enumerate}
		\item Schritt: $x$ Ausklammern
		\item Schritt: Mitternachtsformel in der Klammer anwenden
	\end{enumerate}
	\normalsize\begin{align*}
		2x^3+2x^2-24x&=0\\
		x\cdot(2x^2+2x-24)&=0\\
		=>	x_{1}&= 0\\
		x_{2,3}&= \text{Mitternachtsformel}\\
		=>	x_{2}=3,\;x_{3}&=-4
	\end{align*}	
\end{antwort}
\begin{frage}
	Nullstellen: Funktion 3. Grades
	\[
	f(x)=3x^3+81
	\]
\end{frage}
\begin{antwort}
	Auflösen:
	\begin{align*}
		3x^3+81&=0 \\
		3x^3&=-81 \\
		x^3&=-27 \\
		x&=\sqrt[3]{-27}=-9
	\end{align*}
	\normalsize
	(Die Wurzel einer negativen Zahl ist hier möglich, da es eine ungerade Wurzel ist)
\end{antwort}
\begin{frage}
	Nullstellen: Gebrochene Funktion
	\[
	f(x)=\frac{4x-12}{x^2-4}
	\]
\end{frage}
\begin{antwort}
	Zähler mit 0 gleichsetzen:
	\begin{align*}
	4x-12&=0 \\
	x&=3
	\end{align*}
\end{antwort}
\begin{frage}
	Nullstellen: e-Funktion
	\[
	f(x)=(2x-4)\cdot{e^{3x+4}}
	\]
\end{frage}
\begin{antwort}
	"`e kann niemals 0 werden!"'
	\[
	{e^{3x+4}}\neq0
	\]
	"`Also kann nur der Faktor in der Klammer 0 werden!"'
	\begin{align*}
	2x-4&=0 \\
	x&=2
	\end{align*}
\end{antwort}
\begin{frage}
	Nullstellen: ln-Funktion
	\[
	f(x)=ln(2x-5)
	\]
\end{frage}
\begin{antwort}
	\begin{align*}
		ln(2x-5)&=0\\
		{e^{ln(2x-5)}}&=e^0\\
		2x-5&=1\\
		x&=3
	\end{align*}

\end{antwort}
\end{document}